\topic{solid-ballistics}{Solid-motor internal ballistics: the lumped chamber, its equilibrium, its stability, and how far the chemistry goes}

Reading anchors: Sutton and Biblarz, ninth edition, Section 12.1, printed
pp.~439--457 / PDF pp.~463--481, Eqs.~12--1 to 12--14 \cite{sutton};
Hill and Peterson, second edition, Section 12.7, printed pp.~598--605 / PDF
pp.~610--617, Eqs.~(12.24)--(12.37) \cite{hillpeterson}; Humble,
Section 6.5.1, ``Lumped-Parameter Methods'', printed p.~334 / PDF p.~354
\cite{humble}; Williams, Chapter~7, ``Combustion of Solid Propellants'',
printed pp.~229--261 \cite{williams}.

This topic builds the smallest model of a solid rocket motor that can be
simulated: one chamber, one pressure, one burned-web distance. Everything a
larger model adds --- axial port flow, erosive burning, two-phase nozzle flow,
acoustics, resolved flames --- is a refinement of a term that appears here.
The last subsection answers a question the model raises: how much chemistry
does a solid-motor simulation actually carry?

\subsection{Our notation for the solid motor}
\label{sec:sb-notation}
The subscripts identify the physical role: $p$ in $\rho_p$ denotes
solid propellant, $g$ denotes gas or generation, $c$ denotes the motor
chamber, $b$ denotes the burning surface, and $t$ denotes the throat.
\begin{center}
\begin{tabular}{ll}
Symbol & Meaning\\ \hline
$\rho_p$ & solid propellant density\\
$\rho_g$ & gas density in the motor chamber\\
$V_c$ & gas chamber (cavity) volume\\
$m_c$ & gas mass stored in that chamber\\
$p_c$ & gas chamber pressure (stagnation value for nozzle flow)\\
$A_b$ & burning surface area\\
$r_b$ & normal surface-recession speed\\
$A_t(t)$ & instantaneous nozzle throat area\\
$\dot m_g$ & gas-generation rate at the burning surface\\
$\dot m_{\mathrm{out}}$ & mass-flow rate leaving the chamber\\
$T_c$ & chamber gas temperature (stagnation value; the lumped gas is at rest)\\
$T_{\mathrm{out}}$ & temperature of the gas leaving the control volume\\
$R_s$ & specific gas constant, $k_B/M$ (see below)
\end{tabular}
\end{center}
The motor's gas chamber is the gas-filled space inside its case, including
the cavity exposed as the grain recedes; $V_c$ excludes unburned solid.
A separate chamber downstream of a gas generator is a different control
volume and needs its own subscript, for example
$V_{\mathrm{pl}}$, $p_{\mathrm{pl}}$ and $m_{c,\mathrm{pl}}$.

The symbol $\dot m_g$ denotes the \emph{total} rate at which mass crosses
the burning solid--gas interface into the modelled gas. For uniform
recession and density its local mass flux is $\rho_p r_b$, and its
surface-integrated rate is $\dot m_g=\rho_p A_b r_b$.
The symbol $m_c$ denotes the stored gas inventory: its derivative is
$dm_c/dt$, not $\dot m_g$. The latter is the input to that inventory,
whereas $\dot m_{\mathrm{out}}$ is its output.

We retain $K=A_b/A_t$, initial grain temperature $T_b$, gas heat-capacity
ratio $\gamma$, chamber gas temperature $T_c$, outflow temperature
$T_{\mathrm{out}}$, and specific gas constant $R_s$. In a lumped chamber
the gas leaves at the chamber's stagnation state, so $T_{\mathrm{out}}=T_c$
for its outlets; we keep both symbols because a receiving volume's inflow
temperature is the upstream volume's $T_{\mathrm{out}}$.

\paragraph{Thermodynamic conventions and constants.} We use Kittel and
Kroemer's conventions \cite{CKittelHKroemer1980}, as in Ernest's
thermal-physics notes (\path{mathphysics/LaTeX_and_pdfs/Kittel_Kroemer_Thermal_Physics.tex}
and \path{LaTeXandpdfs/thermo.tex}). The fundamental temperature is
$\tau=k_BT$. A classical ideal gas of $N$ molecules, each of mass $M$, in
volume $V$ obeys
\[
 pV=N\tau,\qquad\text{equivalently}\qquad p=\frac{\rho\tau}{M},
 \qquad \rho=\frac{NM}{V}.
\]
Its heat capacities are $C_V=(\partial U/\partial\tau)_V$ and
$C_p=C_V+N$, dimensionless because $\tau$ is an energy, and
$\gamma=C_p/C_V$; the sound speed satisfies $a^2=\gamma\tau/M$.
The constants in those notes are
\[
 k_B=1.38066\times10^{-23}\ \mathrm{J/K},\quad
 N_0=6.02205\times10^{23}\ \mathrm{mol^{-1}},\quad
 R\equiv N_0k_B=8.31441\ \mathrm{J/(K\,mol)}.
\]
So $R$ without a subscript is the \emph{molar} gas constant. The
2019 SI fixes $k_B=1.380649\times10^{-23}$~J/K and
$N_0=6.02214076\times10^{23}$~mol$^{-1}$ exactly, giving
$R=8.314462618$~J/(K\,mol), the value used in the Cosmos code; the
older values differ in the fifth significant figure.

We keep both temperatures: $\tau$ in energy units, as in those notes,
and $T=\tau/k_B$ in kelvin, as in almost every other physics and
engineering text; so $\tau_c=k_BT_c$ is the chamber gas temperature in
energy units. Time constants are therefore not written $\tau$ here
(Humble's residence time is our $t_R$), and Mach number is written
$\mathfrak M$ (as in \path{LaTeXandpdfs/Propulsion.tex}) because $M$ is
the molecular mass. Define the \emph{specific} gas constant and
specific heats by
\[
 R_s\equiv\frac{k_B}{M},\qquad
 c_v\equiv\frac{k_BC_V}{NM},\qquad
 c_p\equiv\frac{k_BC_p}{NM},
\]
in J/(kg\,K). Then $p=\rho\tau/M=\rho R_sT$, $a^2=\gamma\tau/M=\gamma R_sT$, and $R_sT=\tau/M$,
$\gamma=c_p/c_v$, and $C_p=C_V+N$ becomes
\[
 c_p-c_v=\frac{k_B(C_p-C_V)}{NM}=\frac{k_B}{M}=R_s .
\]
The molar mass is $N_0M$, so $R_s=R/(N_0M)$. Our derivations use these conventions. Textbook equations
retain their authors' symbols and are translated explicitly in
Section~\ref{sec:sb-mass-notation}.


\subsection{Propellant consumption and gas generation from surface recession}
\label{sec:sb-mass-kinematics}
\label{sec:sb-sutton-mass}
Combustion consumes a solid propellant grain as its exposed surface
recedes into the solid. The consumed mass is the density of the solid
times the volume swept out by that moving surface. When all of this
material enters the chamber as gas, conservation of mass identifies the
gas-generation rate with the solid-consumption rate.

We use the following quantities, with their time dependence explicit:
\begin{description}
\item[$S_b(t)$] The exposed burning surface at time $t$.
\item[$A_b(t)$] Area of that burning surface, in $\mathrm m^2$.
\item[$r_b(t)$] Normal surface-recession speed relative to the solid,
in $\mathrm m/\mathrm s$, taken as uniform over $S_b(t)$ initially.
\item[$\rho_p$] Solid propellant density, in $\mathrm{kg}/\mathrm m^3$,
assumed uniform and constant in time.
\item[$m_{\mathrm{burned}}(t)$] Cumulative mass of solid propellant
consumed since a chosen time $t_0$, in $\mathrm{kg}$.
\item[$\dot m_g(t)$] Mass of gas generated per unit time,
in $\mathrm{kg}/\mathrm s$.
\end{description}
The burning area $A_b$ includes only surfaces being consumed; it is
different from the port cross-sectional area and the nozzle throat area.
The recession speed $r_b$ is a length per unit time, whereas
$\dot m_g$ is a mass per unit time. Neither $A_b(t)$ nor
$r_b(t)$ is assumed constant in time.

\paragraph{From recession distance to consumed volume.}
Consider an otherwise stationary solid with a smooth burning surface,
over an interval containing no change of surface topology. In a short
time $\Delta t$, the surface recedes a normal distance
$r_b(t)\Delta t+o(\Delta t)$. Multiplying this distance by the
instantaneous burning area gives, to first order,
\begin{align}
 \Delta V_{\mathrm{consumed}}
 &=A_b(t)r_b(t)\Delta t+o(\Delta t),
 \label{eq:sb-swept-volume}\\
 \Delta m_{\mathrm{burned}}
 &=\rho_p\Delta V_{\mathrm{consumed}}.
 \label{eq:sb-consumed-mass}
\end{align}
Substitute \eqref{eq:sb-swept-volume} into
\eqref{eq:sb-consumed-mass} and divide by $\Delta t$:
\[
 \frac{\Delta m_{\mathrm{burned}}}{\Delta t}
 =\rho_p A_b(t)r_b(t)
   +\rho_p\frac{o(\Delta t)}{\Delta t}.
\]
Taking $\Delta t\to0$ eliminates the last term. If all consumed solid
enters the chamber as gas, with no retained residue or separate surface
storage, conservation of mass then gives
\begin{equation}\label{eq:sb-mass-generation}
 \dot m_g(t)
 =\frac{dm_{\mathrm{burned}}}{dt}
 =\rho_p A_b(t)r_b(t).
\end{equation}
Thus the gas-generation rate is the product of solid density, burning
surface area, and normal recession speed.

\paragraph{When recession is not uniform.}
Let $r_b(\mathbf x,t)$ be the local normal recession speed at
$\mathbf x\in S_b(t)$. Apply the same swept-volume argument to each
surface patch $dA$, then sum over the burning surface:
\begin{equation}\label{eq:sb-surface-mass}
 \frac{dm_{\mathrm{burned}}}{dt}
 =\rho_p\int_{S_b(t)}r_b(\mathbf x,t)\,dA.
\end{equation}
For uniform recession, $r_b$ can be taken outside the integral and
$\int_{S_b(t)}dA=A_b(t)$ recovers
\eqref{eq:sb-mass-generation}. For nonuniform recession, the same product
holds with the area-average
$\bar r_b(t)=A_b(t)^{-1}\int_{S_b(t)}r_b(\mathbf x,t)\,dA$ when $A_b(t)>0$.
If condensed products also enter the chamber, the consumed mass equals
the \emph{total} admitted product mass; obtaining a gas-only generation
rate then requires accounting for the separate phases.

\paragraph{What follows from first principles.}
Equation~\eqref{eq:sb-mass-generation} follows from the definitions of
density, area and recession speed, together with the stated mass-transfer
assumptions. It is a geometric and mass-accounting identity.
Dimensional analysis provides the independent check
\[
 [\rho_p A_b r_b]
 =\frac{\mathrm{kg}}{\mathrm m^3}\,\mathrm m^2\,
   \frac{\mathrm m}{\mathrm s}
 =\frac{\mathrm{kg}}{\mathrm s}.
\]
Dimensions alone would permit a dimensionless multiplier; the
swept-volume derivation fixes that multiplier to one. The identity
does not predict $r_b$. A pressure-dependent burning-rate law supplies
additional physical information, as discussed in
Section~\ref{sec:sb-time-pressure}.

\paragraph{Accumulated burned mass.}
Choose $m_{\mathrm{burned}}(t_0)=0$. Integrating
\eqref{eq:sb-mass-generation} over the interval $[t_0,t]$ gives
\begin{equation}\label{eq:sb-burned-mass-integral}
 m_{\mathrm{burned}}(t)
 =\int_{t_0}^{t}\dot m_g(s)\,ds
 =\rho_p\int_{t_0}^{t}A_b(s)r_b(s)\,ds.
\end{equation}
For a nonzero initial tally, replace the left side by
$m_{\mathrm{burned}}(t)-m_{\mathrm{burned}}(t_0)$.
Only the constant $\rho_p$ has been taken outside the integral.
The remaining solid mass decreases at the same rate:
$dm_{\mathrm{solid}}/dt=-dm_{\mathrm{burned}}/dt$.
This accumulated-mass relation follows from the definition of the rate
and the fundamental theorem of calculus; it adds no combustion law.
Its units are mass, as required. Gas generation and nozzle outflow
can differ during a transient because the chamber can store mass.

\paragraph{Correspondence with the textbooks.}
Hill and Peterson, Section~12.7, printed p.~598 / PDF p.~610,
Eq.~(12.24) \cite{hillpeterson}, present the gas-generation relation with
separate definitions of the generation rate, solid density, burning area
and recession speed. We retain their symbol $\dot m_g$ for surface generation;
their definition of solid density is $\rho_p$ (the inconsistent subscript
printed in Eq.~12.24 is documented in
Section~\ref{sec:sb-mass-crossrefs}).
Sutton and Biblarz, Section~12.1, printed p.~444 / PDF p.~468,
Eqs.~12--1 and 12--2 \cite{sutton}, give the same rate relation and its
time integral, using $\rho_b$ for our $\rho_p$, $\dot m$ for
$\dot m_g$, and $m$ for $m_{\mathrm{burned}}$.
These are parallel presentations of the identities derived here from
surface recession and mass conservation.


\subsection{Time dependence, pressure dependence, and burn history}
\label{sec:sb-time-pressure}
As the grain burns, its exposed area and recession speed can change.
We retain constant solid density while allowing both quantities to vary
with time and pressure, consistently with Sutton, p.~444, and
Hill--Peterson, pp.~598--600 \cite{sutton,hillpeterson}. The instantaneous
notation is explicitly
\[
 A_b=A_b(t),\qquad r_b=r_b(t),\qquad p_c=p_c(t).
\]
The notation $A_b=A_b(t,p_c)$ and $r_b=r_b(t,p_c)$ can be used as
\emph{model shorthand}, evaluated at the gas chamber pressure $p_c(t)$
and with any other state or history understood. Current time and
pressure alone do not generally determine both quantities.

For the rigid grain with uniform regression used below, the more precise
description is
\begin{align}
 r_b(t)&=\mathcal R\bigl(p_c(t),T_b,\ldots\bigr),
 \label{eq:sb-rate-dependence}\\
 y(t)&=y(t_0)+\int_{t_0}^{t}r_b(s)\,ds,
 \label{eq:sb-web-history}\\
 A_b(t)&=\mathcal A\bigl(y(t)\bigr).
 \label{eq:sb-area-history}
\end{align}
Here $\mathcal R$ is a burning-rate law and $\mathcal A$ is a geometry
function; the omitted rate arguments can include local flow or additional
thermal state. The notation $A_b(y)$ below denotes the same geometry
function $\mathcal A(y)$. Substituting
\eqref{eq:sb-rate-dependence} into \eqref{eq:sb-web-history} and then into
\eqref{eq:sb-area-history} makes the pressure-history dependence visible:
\begin{equation}\label{eq:sb-area-pressure-history}
 A_b(t)=\mathcal A\left(
 y(t_0)+\int_{t_0}^{t}\mathcal R\bigl(p_c(s),T_b,\ldots\bigr)\,ds
 \right).
\end{equation}
Two grains with the same initial shape can have the same current pressure
and different burning areas after different pressure histories. In this
rigid model the area has no direct pressure dependence at fixed $y$:
$\left.\frac{\partial A_b}{\partial p_c}\right|_y=0$.
Pressure-dependent deformation would require a different geometry model,
for example $A_b=\mathcal A(y,p_c)$.

For an internally burning cylindrical port of initial radius $R_i$ and
fixed length $L$, with inhibited ends, the geometry gives
\[
 A_b(t)=2\pi L\bigl(R_i+y(t)\bigr),\qquad
 \frac{dA_b}{dt}=2\pi Lr(t)
\]
until the burning surface reaches the outer boundary. This explicitly
shows time dependence through recession. An end-burning grain of constant
cross section can instead have constant $A_b$ while $r_b$ varies.
Substituting \eqref{eq:sb-area-history} into
\eqref{eq:sb-mass-generation} yields the instantaneous rate
$\dot m_g(t)=\rho_p\mathcal A(y(t))r_b(t)$.
There are no $dA_b/dt$ or $dr/dt$ terms in this expression: those arise
only if we differentiate the \emph{rate} again. The geometric identity
does not prescribe $\mathcal R$; $r_b=a(T_b)p_c^n$ below is an additional,
empirical constitutive assumption.

\subsection{Related treatments in the propulsion and combustion texts}
\label{sec:sb-mass-crossrefs}
The following are specific source parallels. The distinction between
definitions, conservation and constitutive assumptions is our assessment
of each passage; it is not a claim that all six books use the same wording.

\paragraph{Hill and Peterson: direct equivalent.}
Section~12.7, printed p.~598 / PDF p.~610, Eq.~(12.24)
\cite{hillpeterson}, equates gas generation with solid consumption and
gives the same density--area--recession product. The supplied scan actually
prints $\rho_g$ in that equation, although the definitions immediately
below name the solid density $\rho_p$ and the subsequent balance on
p.~599 uses $\rho_p$. We use the latter consistently; this is an explicit
notation correction, checked against the source image.
Equation~(12.25), on p.~598, gives the empirical pressure law.
Printed p.~600 / PDF p.~612 explicitly relates changing burning area to
burning rate and initial grain geometry. \emph{Assessment:} the product is
mass accounting; the pressure law is empirical; the changing area follows
from geometry and accumulated recession. These passages do not introduce
a separate counterpart to Sutton's cumulative integral 12--2.

\paragraph{Humble: direct equivalent and explicit history.}
Section~6.5, printed p.~332 / PDF p.~352, Eq.~(6.25)
\cite{humble}, gives
$\dot m_{\mathrm{in}}=\rho_p r_b A_b$ by recalling density times velocity
times flow area. Humble's $r_b$ is also our symbol. The same page explains increasing,
constant and decreasing burning areas using grain geometry.
Section~6.5.1, printed pp.~335--336 / PDF pp.~355--356,
Eq.~(6.31), integrates $r_b(t)$ to obtain web distance and explicitly
connects time-varying burning rate to time-varying pressure.
\emph{Assessment:} Eq.~(6.25) is the same kinematic mass identity;
Eq.~(6.31) integrates a speed to a distance. It supports the history
dependence in \eqref{eq:sb-area-pressure-history}, but its integral is
web distance, not Sutton's accumulated mass.

\paragraph{Huzel and Huang: direct gas-generator equivalent.}
Section~4.6, printed p.~117 / downloaded PDF p.~134
\cite{huzelhuang}, Eq.~(4--46), gives
$\dot w_g=A_bR\rho$, using $R$ for linear burning rate and the book's
\emph{weight flowrate} convention (lb/s, with density in lb/in.$^3$).
For our SI mass balance we use $\dot m_g$, $r_b$, and
$\rho_p$, respectively; a force-valued weight rate would instead be
$g\dot m_g$. Equation~(4--45) specifies pressure dependence
of $R$ at a given propellant temperature. The surrounding discussion
relates gas generation to burning area, pressure and bulk temperature,
and discusses maintaining the area as the grain is consumed.
\emph{Assessment:} Eq.~(4--46) is the same mass-accounting product;
Eq.~(4--45) is a burning-rate correlation. This passage does not give
a separate cumulative-mass integral.

\paragraph{Williams: local flux and regression, not total motor mass.}
Chapter~7, printed p.~230 / PDF p.~253, describes normal surface
regression and the dependence of its speed on pressure, temperature and
composition. Section~7.1, printed p.~232 / PDF p.~255 \cite{williams},
defines the mass burning rate $m$ in $\mathrm{g}/(\mathrm{cm}^2\,\mathrm s)$
in a frame attached to the planar burning interface. Thus his $m$ is a
\emph{mass flux}, unlike Sutton's accumulated mass $m$. In our notation
the translation is $m_{\mathrm{Williams}}=\rho_p r_b$, and integration over
the surface gives \eqref{eq:sb-surface-mass}. Section~7.6,
printed p.~250 / PDF p.~273, Eq.~(7--41), also discusses a
pressure- and temperature-dependent mass-burning-rate expression.
\emph{Assessment:} converting recession to flux is kinematic;
predicting that flux needs combustion physics. The cited planar,
steady passages do not supply a time-dependent whole-grain area or
Sutton's cumulative-mass equation.

\paragraph{Turns: the flux definition and a qualified solid analogue.}
Chapter~3, printed p.~80 / PDF p.~99, Eq.~(3.2) \cite{turns},
defines species mass flux as flowrate per unit area.
Chapter~14, printed p.~534 / PDF p.~553, Eqs.~(14.9)--(14.11),
uses carbon-surface mass accounting and spherical area times flux.
\emph{Assessment:} area times uniform flux is definitional; identifying
the solid's flux with $\rho_p r_b$ uses the recession argument derived
above. This carbon-particle example is not a solid-rocket grain model:
oxygen enters from outside, so outgoing carbon dioxide mass flow is
not the carbon consumption rate. These passages are related
conservation statements, not a presentation of Sutton 12--1/12--2 or
a grain-area law $A_b(t,p_c)$.


\subsection{Chamber mass-flow relation: generation, storage and exhaust}
\label{sec:sb-mass-balance}
Gas entering the chamber from the burning surface has two destinations:
it remains in the chamber or leaves through an outlet. Conservation of
mass relates these rates before any burning-rate law, gas equation of
state or nozzle model is chosen.

In addition to the surface quantities defined in
Section~\ref{sec:sb-mass-kinematics}, define
\begin{description}
\item[$V_c(t)$] Gas chamber (cavity) volume inside the motor, in $\mathrm m^3$.
\item[$\rho_g(t)$] Chamber gas density in the spatially uniform model,
in $\mathrm{kg}/\mathrm m^3$.
\item[$m_c(t)$] Mass of gas currently stored in the chamber, in
$\mathrm{kg}$; this is different from the cumulative burned mass.
\item[$\dot m_{\mathrm{out}}(t)$] Mass leaving this gas chamber per unit time,
in $\mathrm{kg}/\mathrm s$, positive outward.
\item[$p_c(t)$] Pressure of the gas in this motor chamber, in $\mathrm{Pa}$; the stagnation value is used for nozzle flow.
\item[$A_t(t)$] Nozzle throat area, in $\mathrm m^2$.
\item[$c^*(t)$] Characteristic velocity, in $\mathrm m/\mathrm s$.
\end{description}
Retain the single-phase mass-transfer assumptions of
Section~\ref{sec:sb-mass-kinematics}. There is no additional inlet,
leak or retained condensed phase in the present control volume.

\paragraph{Conservation over a time interval.}
Take the control volume to be the gas-filled cavity, whose boundary
moves with the burning surface. Over $[t,t+\Delta t]$,
\[
 m_c(t+\Delta t)-m_c(t)
 =\int_t^{t+\Delta t}
   \bigl[\dot m_g(s)-\dot m_{\mathrm{out}}(s)\bigr]\,ds.
\]
Divide by $\Delta t$ and take $\Delta t\to0$. Rearranging the resulting
rate equation gives
\begin{equation}\label{eq:sb-chamber-continuity}
 \boxed{\dot m_g(t)
 =\frac{dm_c}{dt}+\dot m_{\mathrm{out}}(t).}
\end{equation}
This is generation $=$ accumulation $+$ exhaust. The generation term
denotes mass transferred from solid to gas, not creation of total mass
by a chemical reaction. The fluxes are measured relative to the
control-volume boundary; its recession is already included in the
surface-generation rate.

\paragraph{Stored mass and the lumped approximation.}
For a spatial density field the exact stored mass is
$m_c(t)=\int_{\Omega_g(t)}\rho_g(\mathbf x,t)\,dV$, where
$\Omega_g(t)$ is the gas-filled region. Uniform chamber density reduces
this to
\begin{equation}\label{eq:sb-chamber-mass}
 m_c(t)=\rho_g(t)V_c(t).
\end{equation}
Substitute \eqref{eq:sb-mass-generation} and
\eqref{eq:sb-chamber-mass} into \eqref{eq:sb-chamber-continuity}:
\begin{equation}\label{eq:sb-chamber-product}
 \rho_p A_b(t)r_b(t)
 =\frac{d}{dt}\bigl[\rho_g(t)V_c(t)\bigr]+\dot m_{\mathrm{out}}(t).
\end{equation}
This mass balance does not require constant chamber temperature,
constant pressure, an ideal gas, or a choked outlet.

\paragraph{Keeping an unspecified outlet, or $N$ pintle valves, general.}
Before selecting an outlet-flow model, the relation remains
\[
 \boxed{\rho_p A_b r_b
 =\frac{d(\rho_gV_c)}{dt}+\dot m_{\mathrm{out}}.}
\]
For $N$ outlets on this same chamber, define the total outward rate by
\begin{equation}\label{eq:sb-general-outlets}
 \dot m_{\mathrm{out}}(t)
 \equiv\sum_{i=1}^{N}\dot m_{\mathrm{out},i}(t),
 \qquad
 \dot m_g(t)=\frac{dm_c}{dt}
             +\sum_{i=1}^{N}\dot m_{\mathrm{out},i}(t).
\end{equation}
This sum assumes neither a nozzle geometry nor sonic flow nor a
chemical-composition model. The outlet rates are still unknown.
Conservation relates them to accumulation; it does not determine them
individually or close the transient problem.

A future model might supply a relation of the schematic form
\[
 \dot m_{\mathrm{out},i}
 =\mathcal F_i\bigl(p_{\mathrm{up}},T_{\mathrm{up}},
                   p_{\mathrm{down},i},z_i,\ldots\bigr),
\]
where the upstream pressure and temperature are stagnation quantities,
$p_{\mathrm{down},i}$ is downstream pressure and $z_i$ is actual pintle
position. The ellipsis permits gas composition, losses and additional
flow or actuator states. This is a placeholder for a separately
justified model, not a specified valve law. Measured or prescribed
flow histories could instead supply the rates.

If these valves are attached to a separate receiving chamber, their
upstream state is that chamber's state. The motor's outlet is then
$\dot m_{\mathrm{link}}$, and the two chamber inventories obey
Eq.~\eqref{eq:sb-two-chamber-mass}; the sum of valve rates need not equal
the motor outlet instantaneously. Do not replace an unspecified valve
network by one throat area or insert $p_cA_t/c^*$ without establishing
the corresponding flow model.

\paragraph{Specializing the outlet to an ideal choked nozzle.}
The replacement
\[
 \dot m_{\mathrm{out}}=\frac{p_cA_t}{c^*}
\]
is an additional nozzle relation, not a consequence of chamber mass
conservation. Its elementary predictive derivation uses a
quasi-one-dimensional, quasi-steady, single-phase ideal-gas flow,
with fixed composition and constant heat capacities, adiabatic and
frictionless acceleration to a sonic throat, and no shaft work.
Here $p_c$ is the upstream stagnation pressure, approximated by the
lumped chamber pressure. Section~\ref{sec:sb-pressure-evolution}
derives the formula explicitly from $\dot m=\rho uA$, energy conservation,
the equation of state and isentropy.

\emph{Choked flow} means the throat is sonic:
$\mathfrak M_t=u_t/a_{\mathrm{s},t}=1$, where $u_t$ is local velocity and
$a_{\mathrm{s},t}$ is local sound speed. At fixed upstream state and
throat area, lowering downstream pressure further does not increase
the mass flow while the flow remains choked. It does not mean
$\dot m_{\mathrm{out}}$ is constant in time when upstream conditions
or area change. Choking must be established for the actual boundary
conditions.

\emph{Frozen-composition flow} means species mass fractions remain
unchanged as a gas parcel passes through the nozzle; reactions are
neglected over that transit. It does not freeze pressure, temperature,
velocity, chamber storage or the pintle position. Constant heat
capacities are a further approximation. Frozen chemistry and choking
are independent assumptions: the elementary model here uses both.
An equilibrium-composition or finite-rate-chemistry model may also
have choked flow, but must determine its thermodynamic properties and
mass flux consistently.

Sutton's Section~3.1, assumption~11, printed p.~46 / PDF p.~70
\cite{sutton}, explicitly adopts frozen composition in the nozzle.
Section~5.3, printed p.~167 / PDF p.~191, distinguishes frozen from
shifting-equilibrium composition. The nozzle mass-flow and
characteristic-velocity relations are Eqs.~3--24 and 3--32,
printed pp.~59 and~63 / PDF pp.~83 and~87.

There is also a distinction between a definition and a prediction:
\[
 c^*\equiv\frac{p_cA_t}{\dot m_{\mathrm{out}}}
 \quad\Longleftrightarrow\quad
 \dot m_{\mathrm{out}}=\frac{p_cA_t}{c^*}.
\]
If the first expression defines $c^*$ from measured quantities,
the second is just its rearrangement; it does not establish choking
or supply an unknown flow rate. A prediction requires an independently
specified $c^*$ or flow model valid for the current regime. The ideal
value derived below is one such model, with the stated assumptions.
The notation $p_cA_t/c^*$ alone does not specify a chemistry model.

Only under this nozzle specialization do we substitute into
Eq.~\eqref{eq:sb-chamber-product} to obtain
\begin{equation}\label{eq:sb-motor-mass-flow}
 \boxed{\rho_p A_b(t)r_b(t)
 =\frac{d}{dt}\bigl[\rho_g(t)V_c(t)\bigr]
  +\frac{p_c(t)A_t(t)}{c^*(t)}.}
\end{equation}
The instantaneous notation permits changing throat area and gas
properties. The model below fixes gas properties but retains $A_t(t)$;
only the explicitly frozen-parameter equilibrium and stability analyses
hold the throat area fixed. Freezing parameters in those mathematical
calculations is distinct from freezing the chemical composition
during nozzle transit.
Every term in \eqref{eq:sb-motor-mass-flow} has units of mass per time:
$[p_cA_t/c^*]=\mathrm N/(\mathrm m/\mathrm s)=\mathrm{kg}/\mathrm s$.
Dimensional consistency checks the units; conservation fixes the
equality and the signs.


\paragraph{Storage changes through density and volume.}
Apply the product rule to \eqref{eq:sb-chamber-mass}:
\begin{equation}\label{eq:sb-storage-product-rule}
 \frac{dm_c}{dt}
 =V_c(t)\frac{d\rho_g}{dt}+\rho_g(t)\frac{dV_c}{dt}.
\end{equation}
Thus storage is not caused solely by increasing cavity volume.
Gas can also accumulate by increasing density; a fixed-volume receiving
chamber can store mass. Conversely, a growing cavity does not by itself
guarantee increasing gas mass if its density is falling sufficiently fast.

For a rigid grain and case, with uniform normal recession and no other
change of cavity volume, the swept-volume argument gives
\begin{equation}\label{eq:sb-cavity-growth}
 \frac{dV_c}{dt}=A_b(t)r_b(t).
\end{equation}
Substitute \eqref{eq:sb-cavity-growth} into
\eqref{eq:sb-storage-product-rule}, and then into
\eqref{eq:sb-chamber-product}:
\begin{align}
 \rho_p A_b r_b
 &=V_c\frac{d\rho_g}{dt}+\rho_g A_b r_b+\dot m_{\mathrm{out}},\notag\\
 V_c\frac{d\rho_g}{dt}
 &=(\rho_p-\rho_g)A_b r_b-\dot m_{\mathrm{out}}.
 \label{eq:sb-expanded-mass}
\end{align}
Time arguments have been suppressed only to shorten these two lines.
The factor $\rho_p-\rho_g$ follows by subtracting the gas mass needed
to fill the newly exposed volume at the current gas density. It is
not a change to the solid-generation identity $\dot m_g
=\rho_p A_b r_b$.

\paragraph{Three masses that must remain distinct.}
Integrating \eqref{eq:sb-chamber-continuity} from $t_0$ to $t$ gives
\begin{equation}\label{eq:sb-integrated-chamber-mass}
 m_{\mathrm{burned}}(t)
 =m_c(t)-m_c(t_0)+\int_{t_0}^{t}\dot m_{\mathrm{out}}(s)\,ds,
\end{equation}
where $m_{\mathrm{burned}}(t_0)=0$. Consumed propellant, the change in
stored gas, and expelled mass are separate quantities.
Only when $dm_c/dt=0$ do generation and exhaust balance exactly.
Neglecting storage in a quasi-steady approximation is an additional
approximation, not a consequence of mass conservation. In particular,
$dp_c/dt=0$ alone does not imply $dm_c/dt=0$ in a growing cavity.

For multiple outlets, the same storage argument applies to the total
outlet rate defined in Eq.~\eqref{eq:sb-general-outlets}. For a motor
feeding a receiving chamber, retain both inventories as in
Eq.~\eqref{eq:sb-two-chamber-mass}. No choked-nozzle law is implied
by either conservation statement.

\subsection{The same mass balance in six books: notation and derivations}
\label{sec:sb-mass-notation}
The following notation map concerns the chamber balance. A symbol used
for accumulated burned mass elsewhere in a book need not retain that
meaning here.
\begin{center}
\begin{tabular}{lllll}
Quantity & Ours & Sutton & Hill--Peterson & Humble\\ \hline
solid propellant density & $\rho_p$ & $\rho_b$ & $\rho_p$ & $\rho_p$\\
chamber gas density & $\rho_g$ & $\rho_1$ & $\rho_0$ & $\rho_g$\\
gas chamber (cavity) volume & $V_c$ & $V_1$ & $\mathcal V_0$ & $V$\\
stored chamber gas mass & $m_c$ & $\rho_1V_1$ & $\mathcal M_s$ & $m$\\
gas chamber pressure & $p_c$ & $p_1$ & $p_0$ & $p_c$\\
recession speed & $r_b$ & $r$ & $r$ & $r_b$\\
throat area & $A_t$ & $A_t$ & $A^*$ & $A_t$\\
surface generation rate & $\dot m_g$ & $\rho_b A_b r$
 & $\dot m_g$ & $\dot m_{\mathrm{in}}$\\
chamber outlet rate & $\dot m_{\mathrm{out}}$ & $p_1A_t/c^*$
 & $\dot m_n$ & $\dot m_{\mathrm{out}}$
\end{tabular}
\end{center}

\paragraph{Hill and Peterson: expanding the chamber storage.}
Section~12.7, printed p.~599 / PDF p.~611 \cite{hillpeterson},
first identifies the stored gas mass and the conservation balance:
\[
 \mathcal M_s=\rho_0\mathcal V_0,\qquad
 \dot m_g=\frac{d\mathcal M_s}{dt}+\dot m_n.
\]
Using $d\mathcal V_0/dt=rA_b$, the book then writes
\[
 \rho_p A_b r
 =\rho_0 A_b r+\mathcal V_0\frac{d\rho_0}{dt}+\dot m_n.
\]
These are \eqref{eq:sb-chamber-continuity} and
\eqref{eq:sb-expanded-mass} before rearrangement. Substitution of the
empirical law $r=ap_0^n$ and the choked-nozzle flow from Eq.~(3.14)
produces their Eq.~(12.27):
\begin{align*}
 \rho_p A_b a p_0^n
 &=\rho_0 A_b a p_0^n+\mathcal V_0\frac{d\rho_0}{dt}\\
 &\quad+\frac{p_0 A^*}{\sqrt{RT_0}}\sqrt{\gamma}
 \left(\frac{2}{\gamma+1}\right)^{\frac{\gamma+1}{2(\gamma-1)}}.
\end{align*}
Their next step uses fixed $T_0$ and $R$ (our $T_c$ and $R_s$) to replace
$d\rho_0/dt$ by $(RT_0)^{-1}dp_0/dt$, giving Eq.~(12.28).
Thus the conservation equation precedes the burning and nozzle closures.

\paragraph{Humble: inflow minus outflow.}
Section~6.5.1, printed p.~334 / PDF p.~354 \cite{humble},
Eqs.~(6.26)--(6.27), writes
\[
 \dot m_{\mathrm{out}}=\frac{p_c A_t}{c^*},\qquad
 \frac{dm}{dt}
 =\dot m_{\mathrm{in}}-\dot m_{\mathrm{out}}
 =r_b\rho_p A_b-\frac{p_c A_t}{c^*}.
\]
Here $m$ is stored chamber gas, not consumed propellant.
The unsteady derivation on printed p.~338 / PDF p.~358,
Eqs.~(6.33)--(6.36), starts from
\begin{align*}
 m&=\frac{p_c V\mathcal M}{\mathcal R_uT_c},\\
 \frac1m\frac{dm}{dt}
 &=\frac1{p_c}\frac{dp_c}{dt}+\frac1V\frac{dV}{dt},\\
 \frac{dV}{dt}&=r_bA_b.
\end{align*}
Humble's $\mathcal M$ is the molar mass in kg/kmol, $10^3N_0M$ in our
notation, and his $\mathcal R_u=8314$~J/(kmol\,K) is $10^3R$. Hence
$\mathcal R_u/\mathcal M=R/(N_0M)=k_B/M=R_s$, and his first line is
$m_c=p_cV_c/(R_sT_c)$. His $T_c$ is our $T_c$. The logarithmic
derivative shown assumes fixed $T_c$ and fixed gas composition.
Solving for the pressure derivative and inserting the mass balance gives
\begin{align*}
 \frac1{p_c}\frac{dp_c}{dt}
 &=\frac1m\left(r_b\rho_p A_b-\frac{p_cA_t}{c^*}\right)
   -\frac{r_bA_b}{V}\\
 &=\frac1m\left[r_bA_b(\rho_p-\rho_g)
                  -\frac{p_cA_t}{c^*}\right],
\end{align*}
where $\rho_g=m/V$. The last line is Eq.~(6.36), the same
density-difference correction obtained in \eqref{eq:sb-expanded-mass}.

\paragraph{Sutton and Biblarz: substituting each physical term.}
Section~12.1, printed pp.~444--445 / PDF pp.~468--469
\cite{sutton}, Eq.~12--3, expresses the same balance as
\eqref{eq:sb-chamber-continuity}. Begin with our rate equation, before
substituting any component formula:
\begin{equation}\label{eq:sb-sutton-start}
 \underbrace{\dot m_g}_{\text{surface generation}}
 =\underbrace{\frac{dm_c}{dt}}_{\text{gas accumulation}}
  +\underbrace{\dot m_{\mathrm{out}}}_{\text{outlet flow}}.
\end{equation}
The middle term is the rate of change of \emph{stored gas mass}.
It can be positive, zero or negative. It is not the surface-generation
rate $\dot m_g$.

The first two substitutions define generation and inventory under the
stated lumped assumptions. The third is the additional nozzle model
needed to reproduce Sutton's specialized expression:
\begin{align}
 \dot m_g&=\rho_p A_b r_b
 &&\text{surface recession, Eq.~\eqref{eq:sb-mass-generation}},
 \label{eq:sb-sutton-generation}\\
 \frac{dm_c}{dt}
 &=\frac{d}{dt}(\rho_gV_c)
 &&\text{stored mass, Eq.~\eqref{eq:sb-chamber-mass}},
 \label{eq:sb-sutton-storage}\\
 \dot m_{\mathrm{out}}&=\frac{p_c A_t}{c^*}
 &&\text{choked nozzle, Eq.~\eqref{eq:sb-nozzle-flow}}.
 \label{eq:sb-sutton-outlet}
\end{align}
The first formula uses uniform surface recession and constant solid
density; the second uses spatially uniform chamber gas density.
The third formula uses the ideal nozzle model: \emph{choked} means
sonic flow at the throat, and \emph{frozen composition} means species
fractions do not change during nozzle transit. The derivation in
Section~\ref{sec:sb-pressure-evolution} additionally assumes quasi-steady,
adiabatic, frictionless flow with constant gas heat capacities.
For unspecified pintle outlets, keep $\dot m_{\mathrm{out}}$ and stop
before the final substitution below:
\begin{align*}
 \dot m_g&=\frac{dm_c}{dt}+\dot m_{\mathrm{out}}\\
 \rho_p A_b r_b&=\frac{dm_c}{dt}+\dot m_{\mathrm{out}}
 &&\text{insert \eqref{eq:sb-sutton-generation}},\\
 \rho_p A_b r_b&=\frac{d(\rho_gV_c)}{dt}+\dot m_{\mathrm{out}}
 &&\text{insert \eqref{eq:sb-sutton-storage}},\\
 \rho_p A_b r_b&=\frac{d(\rho_gV_c)}{dt}+\frac{p_c A_t}{c^*}
 &&\text{insert \eqref{eq:sb-sutton-outlet}}.
\end{align*}
The third line is the general chamber balance; the last line is its
nozzle specialization, Eq.~\eqref{eq:sb-motor-mass-flow}, with the time
arguments suppressed. For that nozzle model, written out explicitly, it is
\[
 \rho_p A_b(t)r_b(t)
 =\frac{d}{dt}\bigl[\rho_g(t)V_c(t)\bigr]
  +\frac{p_c(t)A_t(t)}{c^*(t)}.
\]

Sutton's symbols map to ours as follows:
\[
 \underbrace{\rho_b}_{\text{Sutton}}\longleftrightarrow
 \underbrace{\rho_p}_{\text{ours}},\qquad
 \rho_1\longleftrightarrow\rho_g,\qquad
 V_1\longleftrightarrow V_c,\qquad
 p_1\longleftrightarrow p_c.
\]
The symbols $A_b,A_t,c^*$ are unchanged, and Sutton's $r$ is our
$r_b$. The two notations can
therefore be read side by side:
\begin{center}
\begin{tabular}{c@{\qquad}c}
Sutton, Eq.~12--3 & Our notation\\[6pt]
$\displaystyle A_b r\rho_b
 =\frac{d(\rho_1V_1)}{dt}+\frac{A_t p_1}{c^*}$
&
$\displaystyle \rho_p A_b r_b
 =\frac{d(\rho_gV_c)}{dt}+\frac{p_c A_t}{c^*}$
\end{tabular}
\end{center}
In particular, the storage term expands mathematically as
\[
 \frac{dm_c}{dt}
 =\frac{d(\rho_gV_c)}{dt}
 =V_c\frac{d\rho_g}{dt}+\rho_g\frac{dV_c}{dt}
 =V_c\frac{d\rho_g}{dt}+\rho_g A_b r_b,
\]
where the last equality uses \eqref{eq:sb-cavity-growth}. Thus the
same substitution chain also gives
\[
\begin{aligned}
 \rho_p A_b r_b
 &=V_c\frac{d\rho_g}{dt}+\rho_g A_b r_b+\frac{p_c A_t}{c^*},\\
 V_c\frac{d\rho_g}{dt}
 &=(\rho_p-\rho_g)A_b r_b-\frac{p_c A_t}{c^*}.
\end{aligned}
\]
This shows explicitly which part of the generated mass changes gas
density, which part fills the growing gas chamber, and which part
leaves through the nozzle.


\paragraph{Huzel and Huang: stored inventory in an engine model.}
Section~10.2, printed p.~347 / downloaded PDF p.~361
\cite{huzelhuang}, Eqs.~(10--7)--(10--10), stores oxidizer-origin and
fuel-origin inventories $w_O,w_f$, with $w_C=w_O+w_f$ and
$R_O=w_O/w_C$. After suppressing the injector transport delays, their
inlet and outlet bookkeeping is equivalently
\begin{align*}
 \frac{dw_O}{dt}&=\dot w_O-R_O\dot w_C,\\
 \frac{dw_f}{dt}&=\dot w_f-(1-R_O)\dot w_C.
\end{align*}
Adding the two equations and using $R_O+(1-R_O)=1$ gives
\[
 \frac{dw_C}{dt}=\dot w_O+\dot w_f-\dot w_C.
\]
This is our differentiated, delay-free form of their inventory
integrals. In this source notation, $w_C$ denotes stored inventory,
whereas $\dot w_C$ denotes \emph{outlet flow}; it is not shorthand for
$dw_C/dt$. The quantities use the book's weight/flow convention
(lb and lb/s); translating them consistently to SI mass quantities gives
$dm_c/dt=\dot m_{\mathrm{in}}-\dot m_{\mathrm{out}}$.
For a solid gas generator, Section~4.6, p.~117, Eq.~(4--46), supplies
the inlet counterpart $\dot w_g=A_bR\rho$, where this $R$ is recession
speed. Replacing the two liquid-feed terms by this single surface feed
is our adaptation of the inventory balance to a solid motor:
\[
 \dot w_g\longleftrightarrow\dot m_g,\qquad
 \dot w_C\longleftrightarrow\dot m_{\mathrm{out}},\qquad
 \frac{dm_c}{dt}=\dot m_g-\dot m_{\mathrm{out}}.
\]

Their Eq.~(10--11) intends the nozzle outlet relation
$\dot w_C=A_tP_Cg/C^*$ in their unit convention. The supplied page
prints $w_C$ without a dot, but the flow definition, dimensions and
Eq.~(1--32) identify an outlet rate; this correction is explicit here.
The conversion factor $g$ in that expression must not be inserted
into the SI mass-flow formula $p_cA_t/c^*$.

\paragraph{Turns: summing species balances to conserve total mass.}
Chapter~6, printed pp.~194--195 / PDF pp.~213--214
\cite{turns}, Eq.~(6.27), writes for species $i$
\[
 \frac{dm_{i,cv}}{dt}
 =\dot m_i'''V+\dot m_{i,\mathrm{in}}-\dot m_{i,\mathrm{out}}.
\]
Here $m_{i,cv}$ is stored species mass and $\dot m_i'''$ is its
volumetric chemical production rate. Chemical reactions redistribute
mass among species, so $\sum_i\dot m_i'''=0$. Summing the species
equations and identifying $m_{cv}=\sum_i m_{i,cv}$ therefore gives
Eq.~(6.28):
\[
 \frac{dm_{cv}}{dt}=\dot m_{\mathrm{in}}-\dot m_{\mathrm{out}}.
\]
The explicit map to our solid-motor balance is
\[
 m_{cv}\longleftrightarrow m_c=\rho_gV_c,\qquad
 \dot m_{\mathrm{in}}\longleftrightarrow\dot m_g,\qquad
 \frac{dm_c}{dt}=\dot m_g-\dot m_{\mathrm{out}}.
\]
Turns' $\dot m_{\mathrm{out}}$ already agrees with our outlet notation. Applying it to a receding
surface requires boundary-relative fluxes; a gas-phase reaction source
must not be counted again as additional total mass generation.

\paragraph{Williams: the moving-control-volume derivation.}
Section~1.1, printed p.~2 / PDF p.~25 \cite{williams},
Eq.~(1--1), gives
$\partial\rho/\partial t+\nabla\cdot(\rho\mathbf v)=0$.
Section~1.4, printed p.~13 / PDF p.~36, Eqs.~(1--50)--(1--51),
integrates this continuity equation over a moving volume and applies
the transport theorem and divergence theorem:
\[
 \frac{d}{dt}\int_{\mathcal V}\rho\,d\mathcal V
 +\int_{\mathcal A}\rho
   \left(\mathbf v-\frac{d\mathbf x}{dt}\right)
   \cdot\mathbf n\,d\mathcal A=0.
\]
Williams uses $\mathcal V$ for the control volume, $\mathcal A$ for its
bounding surface, $\mathbf v$ for mass-average gas velocity,
$d\mathbf x/dt$ for boundary velocity, and $\mathbf n$ for the outward
normal. His $\mathcal A$ here is not the scalar burning area $A_b$.
For the motor cavity the volume integral is $m_c$, and the boundary
flux integral splits into
\[
 \int_{\mathcal A}\rho
   \left(\mathbf v-\frac{d\mathbf x}{dt}\right)\cdot\mathbf n\,d\mathcal A
 =\dot m_{\mathrm{out}}-\dot m_g.
\]
The fixed impermeable walls contribute zero; the burning surface is
an inlet and the nozzle is an outlet. Substitution recovers
\eqref{eq:sb-chamber-continuity}, with its signs fixed by the outward
normal. This is the general conservation foundation of the lumped
motor equation, specialized here to our control volume.


\subsection{From transient storage to quasi-steady burning: what may be neglected}
\label{sec:sb-storage-approximation}
Sutton's transition from Eq.~12--3 to Eq.~12--4, printed
p.~445 / PDF p.~469 \cite{sutton}, makes an approximation to the
mass balance. It does not introduce a new conservation law.
The preceding discussion explicitly identifies startup and shutdown as
transient applications of Eq.~12--3 and says storage matters during
startup. The later short-duration qualification should be read in that
context; duration alone is not a sufficient test for neglecting storage.

\paragraph{Which part does the density ratio make small?}
For the same small mass $\Delta m$, compare its solid and gas volumes
at the respective densities:
\[
 \Delta V_p=\frac{\Delta m}{\rho_p},\qquad
 \Delta V_g=\frac{\Delta m}{\rho_g},\qquad
 \frac{\Delta V_g}{\Delta V_p}=\frac{\rho_p}{\rho_g}.
\]
Thus Sutton's approximate thousandfold volume expansion is a comparison
between solid and gas densities at the conditions considered, not a
claim about an initially small number of molecules in an empty chamber.
The ratio is not universal: $\rho_g$ depends on the gas state.

Using the product rule and the recession relation gives
\begin{align}
 \frac{dm_c}{dt}
 &=\underbrace{V_c\frac{d\rho_g}{dt}}_{\text{density change}}
   +\underbrace{\rho_g\frac{dV_c}{dt}}_{\text{volume change}}
   =V_c\frac{d\rho_g}{dt}+\rho_g A_b r_b,
   \label{eq:sb-storage-two-terms}\\
 \frac{\rho_g\,dV_c/dt}{\dot m_g}
 &=\frac{\rho_g A_b r_b}{\rho_p A_b r_b}
   =\frac{\rho_g}{\rho_p}\ll1,
   \qquad A_b r_b>0.
   \label{eq:sb-volume-storage-ratio}
\end{align}
The small density ratio controls the volume-change term only.
It places no bound on $V_c\,d\rho_g/dt$. In particular,
\[
 \rho_g\ll\rho_p
 \quad\not\Longrightarrow\quad
 \left|\frac{d(\rho_gV_c)}{dt}\right|\ll\dot m_g.
\]
For example, while generation is positive but outlet flow is very small,
the exact mass balance gives
\[
 \dot m_{\mathrm{out}}\ll\dot m_g
 \quad\Longrightarrow\quad
 \frac{dm_c}{dt}
 =\dot m_g-\dot m_{\mathrm{out}}\simeq\dot m_g.
\]
Storage is then a leading term even if the instantaneous gas density is
low. This can occur during initial pressurization or outlet closure.
A small inventory and a small rate of change of inventory are different
statements.

\paragraph{A quantitative test for negligible storage.}
During active burning, $\dot m_g>0$, define
\begin{align}
 \epsilon_{\mathrm{store}}
 &\equiv\frac{|dm_c/dt|}{\dot m_g}
  =\left|1-\frac{\dot m_{\mathrm{out}}}{\dot m_g}\right|,
  \label{eq:sb-storage-test}\\
 \epsilon_\rho
 &\equiv\frac{V_c|d\rho_g/dt|}{\dot m_g},\qquad
 \epsilon_V\equiv\frac{\rho_g A_b r_b}{\dot m_g}
  =\frac{\rho_g}{\rho_p},\notag\\
 \epsilon_{\mathrm{store}}&\le\epsilon_\rho+\epsilon_V.\notag
\end{align}
Neglecting accumulation requires $\epsilon_{\mathrm{store}}\ll1$;
making both contributions individually small is a sufficient test.
At burnout $\dot m_g=0$, this normalization is undefined and the
unscaled balance must be retained. Gas can continue to leave through
the outlet by drawing down the stored inventory.

\paragraph{Eq.~12--4: negligible accumulation and a choked nozzle.}
Here \emph{choked} means sonic flow at the controlling throat:
$\mathfrak M_t=u_t/a_{\mathrm{s},t}=1$. With upstream pressure, temperature,
composition and throat area fixed, lowering downstream pressure further
does not increase mass flow while choking persists. Frozen chemistry,
which holds composition fixed during transit, is a separate assumption
in the elementary $c^*$ model (Section~\ref{sec:sb-mass-balance}).
Start with the full instantaneous balance and rearrange:
\begin{align}
 \dot m_g(t)
 &=\frac{dm_c}{dt}+\dot m_{\mathrm{out}}(t),\notag\\
 \rho_p A_b(t)r_b(t)
 &=\frac{dm_c}{dt}+\frac{p_c(t)A_t(t)}{c^*(t)},\notag\\
 p_c(t)
 &=K(t)\rho_p r_b(t)c^*(t)
   -\frac{c^*(t)}{A_t(t)}\frac{dm_c}{dt},
 \quad K(t)\equiv\frac{A_b(t)}{A_t(t)}.
 \label{eq:sb-pressure-with-storage}
\end{align}
Here $A_t(t)>0$. To reach Eq.~12--4, make the further approximation
that accumulation is negligible compared with generation and outflow.
The two distinct assumptions are then both in force:
\[
 \underbrace{\dot m_c\equiv\frac{dm_c}{dt}\simeq0}_{
                    \text{negligible accumulation}},
 \qquad
 \underbrace{\dot m_{\mathrm{out}}=\frac{p_cA_t}{c^*}}_{
                    \text{choked-nozzle model}}.
\]
The first is an approximation on a rate, not a statement that
$m_c=0$. The second supplies the outlet model; the elementary
thermodynamic value of $c^*$ additionally uses the frozen-composition
and constant-heat-capacity assumptions stated earlier.
Substituting both into the mass balance gives, step by step,
\[
 \dot m_g\simeq\dot m_{\mathrm{out}}
 \quad\Longrightarrow\quad
 \rho_pA_b r_b\simeq\frac{p_cA_t}{c^*}
 \quad\Longrightarrow\quad
 p_c\simeq\frac{A_b}{A_t}\rho_p r_bc^*.
\]
Thus
\begin{equation}\label{eq:sb-sutton-quasisteady}
 \boxed{p_c(t)\simeq K(t)\rho_p r_b(t)c^*(t).}
\end{equation}
This is Sutton's Eq.~12--4 in our notation: a specialized approximation,
not the general motor or valve-network balance. The approximation sign
and time dependence are explicit. No step required $dA_t/dt=0$.
It is generally an implicit relation because $r_b$ depends on $p_c$.
If a burning law is supplied, the pressure given by this algebraic
balance is a quasi-steady target, not automatically the transient
pressure.

\paragraph{Constant pressure, negligible storage, and constant area differ.}
With fixed gas temperature and composition, constant $p_c$ implies
constant $\rho_g$. Nevertheless, a burning grain increases $V_c$, so
\[
 \frac{dp_c}{dt}=0
 \ \Longrightarrow\
 \frac{d\rho_g}{dt}=0,\qquad
 \frac{dm_c}{dt}=\rho_g A_b r_b>0.
\]
Thus even constant-pressure burning has a small positive storage term
when the cavity grows. Dropping it is an approximation justified by
$\rho_g/\rho_p\ll1$, not an exact consequence of constant pressure.
During a slowly evolving burn, $A_b(t)$ and $A_t(t)$ can change while
the pressure follows a sequence of approximate equilibria.

\paragraph{What $K(t)$ means when either area changes.}
Sutton's definition, Eq.~12--4 on printed p.~445 / PDF p.~469
\cite{sutton}, is a dimensionless area ratio. Its instantaneous form is
\begin{equation}\label{eq:sb-k-variable}
 \boxed{K(t)\equiv\frac{A_b(t)}{A_t(t)},\qquad A_t(t)>0.}
\end{equation}
This definition requires neither constant throat area nor steady
burning nor negligible storage nor choked flow. Both areas can vary.
It is naturally useful for a surface-burning gas generator, including
a solid motor with a variable controlling throat. Its definition is
geometric; using it to predict pressure requires the additional
generation and outlet models.

Apply the quotient rule:
\begin{align}
 \frac{dK}{dt}
 &=\frac{A_t\,dA_b/dt-A_b\,dA_t/dt}{A_t^2}\notag\\
 &=\frac{1}{A_t}\frac{dA_b}{dt}
   -\frac{A_b}{A_t^2}\frac{dA_t}{dt}.
 \label{eq:sb-k-rate}
\end{align}
When $A_b>0$ as well, divide by $K=A_b/A_t$:
\[
 \frac{1}{K}\frac{dK}{dt}
 =\frac{1}{A_b}\frac{dA_b}{dt}
  -\frac{1}{A_t}\frac{dA_t}{dt}.
\]
For example, comparing two instants with the same burning area,
\[
 A_b(t_2)=A_b(t_1),\quad A_t(t_2)=2A_t(t_1)
 \quad\Longrightarrow\quad K(t_2)=\frac{K(t_1)}{2}.
\]
This is an area-ratio statement; it does not assert that the transient
pressure immediately halves. Constant $K$ requires equal fractional
rates of change of the two areas, not that either area be constant.
Sutton's parenthetical association of constant $K$ with constant $A_b$
uses the preliminary fixed-throat assumption.

The phrase \emph{at fixed throat area} in the definition of choking
holds geometry fixed while comparing the effect of different downstream
pressures. It does not impose $dA_t/dt=0$ throughout a firing. A
variable-area nozzle can remain choked as its area and mass flow change.
Likewise, a well-defined $K(t)$ does not establish choking or justify
Eq.~\eqref{eq:sb-sutton-quasisteady}.
Inferring a nearly constant pressure from nearly constant $K$ also
requires a suitable burning law, control of the gas-property parameters,
and negligible accumulation. These are additional model assumptions.
Steady burning in that approximation concerns an operating regime,
not the definition of $K$.

\paragraph{Several outlets: a geometric ratio and its conditional use.}
For $N$ parallel outlets directly on the same lumped motor chamber,
one can define the total geometric throat area and its corresponding
ratio:
\begin{equation}\label{eq:sb-k-multiple}
 A_{t,\Sigma}(t)\equiv\sum_{i=1}^{N}A_{t,i}(t),\qquad
 K_\Sigma(t)\equiv\frac{A_b(t)}{A_{t,\Sigma}(t)},
 \quad A_{t,\Sigma}(t)>0.
\end{equation}
The usefulness of this ratio for the flow follows under specific
assumptions. Suppose every active outlet is choked (sonic at its
controlling throat), each is described by the ideal quasi-steady
nozzle model, and all share the same upstream stagnation pressure,
temperature and gas composition. Their ideal $c^*(t)$ is then common.
Using Eq.~\eqref{eq:sb-nozzle-flow} for each outlet and summing gives
\begin{align}
 \dot m_{\mathrm{out},i}(t)
 &=\frac{p_c(t)A_{t,i}(t)}{c^*(t)},\notag\\
 \dot m_{\mathrm{out}}(t)
 &=\sum_{i=1}^{N}\frac{p_c(t)A_{t,i}(t)}{c^*(t)}
  =\frac{p_c(t)}{c^*(t)}A_{t,\Sigma}(t),\notag\\
 \rho_pA_b r_b
 &=\frac{dm_c}{dt}+\frac{p_cA_{t,\Sigma}}{c^*}.
 \label{eq:sb-multiple-choked-balance}
\end{align}
This is our direct generalization of the single-nozzle balance.
It retains accumulation and allows the individual areas to vary.
Only with negligible accumulation does it reduce to
\[
 p_c(t)\simeq K_\Sigma(t)\rho_p r_b(t)c^*(t).
\]
Thus a single nozzle is not essential, but common upstream conditions
and compatible outlet laws are essential to this particular reduction.

If separate discharge corrections $C_{d,i}(t)$ are specified relative
to the same ideal $c^*(t)$, and each outlet is modelled by
$\dot m_{\mathrm{out},i}=C_{d,i}p_cA_{t,i}/c^*$, then the corresponding
definitions are instead
\[
 A_{\mathrm{eff}}(t)
 \equiv\sum_{i=1}^{N}C_{d,i}(t)A_{t,i}(t),\qquad
 K_{\mathrm{eff}}(t)\equiv\frac{A_b(t)}{A_{\mathrm{eff}}(t)}.
\]
Here $A_{\mathrm{eff}}>0$, and $C_{d,i}$ is a dimensionless correction
to that ideal mass flow. This effective ratio is model-dependent,
not solely geometric. It does not supply the correction factors,
the pintle-position-to-area map, or proof of choking.

\paragraph{Where a single $K$ stops being sufficient.}
An unchoked outlet, unequal upstream states or an unspecified valve
law prevents the common factorization above from being assumed.
A sum of areas may still be a geometric descriptor, but is not then
a demonstrated equivalent nozzle. Retain
\[
 \dot m_g=\frac{dm_c}{dt}+\sum_{i=1}^{N}\dot m_{\mathrm{out},i}
\]
until appropriate models or data supply the individual rates.
For a motor feeding a separate receiving chamber, its pressure need
not equal the motor pressure, and its valve outflow need not equal
the link flow during a transient. A ratio formed with motor burning
area and downstream valve areas cannot remove those extra states
and pressure drops; use Eq.~\eqref{eq:sb-two-chamber-mass}.

If the chosen denominator area is zero, the corresponding $K$ is
undefined. This is a limit of the ratio, not a prediction of infinite
physical pressure. The undivided mass balance remains valid.
Even when defined, $K$ alone does not encode chamber volume, stored
mass, gas properties or valve losses, so it cannot specify the full
transient response.


\paragraph{Assessment and scope for a pintle-controlled system.}
Sutton's Eq.~12--4 is a valid leading approximation under the stated
small-storage and nozzle assumptions. The density-ratio argument
alone does not justify neglecting the entire derivative, and the
short-operating-duration exception is incomplete as a general
criterion: startup, shutdown, or a sufficiently rapid valve movement
can make accumulation important during a long firing.
Our transient equations therefore retain storage and allow time-varying
outlet area. Fixed-area calculations below are explicitly identified
as frozen-parameter analyses.

The geometric nozzle throat $A_t(t)$ must also be distinguished from
a valve's effective flow area. A pintle can alter the controlling
restriction, but its displacement need not equal an area change
linearly, and losses need their own flow model. The formula
$p_cA_t/c^*$ is not automatically the flow law of every valve or of
an unchoked connection. A downstream receiving chamber has its own
pressure and gas inventory; its valves respond to that chamber's
state. Their mass balances, with $\dot m_{\mathrm{link}}$ denoting
flow from motor to receiving chamber, are
\begin{align}
 \frac{dm_c}{dt}
 &=\dot m_g-\dot m_{\mathrm{link}},\notag\\
 \frac{dm_{c,\mathrm{pl}}}{dt}
 &=\dot m_{\mathrm{link}}
   -\sum_{i=1}^{N}\dot m_{\mathrm{out},i},\notag\\
 \frac{d(m_c+m_{c,\mathrm{pl}})}{dt}
 &=\dot m_g-\sum_{i=1}^{N}\dot m_{\mathrm{out},i}.
 \label{eq:sb-two-chamber-mass}
\end{align}
The link flow cancels when the two inventories are added. Neither
inventory may be discarded solely because the solid-to-gas density
ratio is large.


\subsection{State, assumptions and pressure evolution}
\label{sec:sb-pressure-evolution}
The state is $(p_c,y)\in(0,\infty)\times[0,w]$, where $y$ is the distance
the burning surface has receded normal to itself and $w$ is the web
thickness. Initially keep $\dot m_{\mathrm{out}}$ as an unspecified rate.
For the nozzle specialization, the area $A_t(t)$ is a prescribed input;
predicting it from erosion or actuator motion requires an additional model.
The grain geometry is given as two functions of $y$ alone:
$A_b(y)>0$ and $V_c(y)>0$. Their compatibility condition is
\begin{equation}\label{eq:sb-volume-rate}
 \frac{dV_c}{dy}=A_b(y),
\end{equation}
which says that each unit of recession uncovers a slab of volume $A_b\,dy$.

Assumptions:
\begin{itemize}
\item[B1] \emph{Lumped chamber.} $p_c$, $\rho_g$ and $T_c$ are uniform in the
cavity. (Hill--Peterson's axial pressure drop, treated below, is the first
correction.)
\item[B2] \emph{Perfect gas with fixed $T_c$, $R_s$, $\gamma$.} Hence
$\rho_g=p_c/(R_sT_c)$. Hill and Peterson justify fixing it (their $T_0$): it ``is
practically independent of changes in combustion pressure for typical solid
propellants'' (p.~599).
\item[B3] \emph{Optional nozzle model: choked, quasi-steady flow with area $A_t(t)>0$.}
At the throat $\mathfrak M_t=1$ (choked flow); composition is unchanged during
transit (frozen chemistry). For the ideal formula below the gas has
constant heat capacities and accelerates adiabatically without friction
to that throat. The nozzle adjusts sufficiently quickly for an instantaneous
flow law to apply; the chamber inventory still evolves dynamically. Area is not
assumed constant. Erosion or a moving restriction can change it, provided
this nozzle description remains appropriate. Sutton's reported erosion
range on p.~445 is not a bound on a pintle's commanded area variation.
\item[B4] \emph{Saint-Robert (Vieille) burning law}, Sutton Eq.~12--5,
Hill--Peterson Eq.~(12.25):
\begin{equation}\label{eq:sb-burn-law}
 r_b\equiv\frac{dy}{dt}=a(T_b)\,p_c^n,\qquad a>0,\ n\ \text{constant}.
\end{equation}
No erosive burning: $r_b$ does not depend on the gas velocity over the surface.
\item[B5] $\rho_p$ constant.
\end{itemize}

\paragraph{Nozzle flow: deriving the optional outlet relation.}
This paragraph uses B3 in addition to the chamber assumptions.
For the elementary nozzle calculation take a single-phase gas with
fixed composition, constant specific heat $c_p$, and constant $R_s$ and
$\gamma>1$, with
$c_p=\gamma R_s/(\gamma-1)$. Use adiabatic, frictionless,
quasi-one-dimensional flow from the chamber to a sonic throat, with
no shaft work. The quasi-steady approximation means this passage
adjusts fast enough to apply the steady equations at each time;
it does not set $dm_c/dt$ to zero.

Let $p_t,T_t,\rho_t,u_t$ be the local static pressure, temperature,
density and velocity at the throat. The inlet stagnation quantities
are $p_c,T_c$. In the notation of Section~\ref{topic:nozzle},
$p_0\leftrightarrow p_c$ and $T_0\leftrightarrow T_c$; $R_s$ is
the same symbol.
For this gas the throat sound speed is
$a_{\mathrm{s},t}=\sqrt{\gamma R_sT_t}$. Choking imposes
\begin{equation}\label{eq:sb-sonic-throat}
 \mathfrak M_t=\frac{u_t}{a_{\mathrm{s},t}}=1,
 \qquad u_t=\sqrt{\gamma R_sT_t}.
\end{equation}
Now substitute this speed into the energy equation:
\begin{align*}
 c_pT_c&=c_pT_t+\frac{u_t^2}{2}
        =c_pT_t+\frac{\gamma R_sT_t}{2},\\
 T_c&=T_t\left(1+\frac{\gamma R_s}{2c_p}\right)
     =T_t\frac{\gamma+1}{2},\\
 \frac{T_t}{T_c}&=\frac{2}{\gamma+1},\qquad
 \frac{p_t}{p_c}
 =\left(\frac{T_t}{T_c}\right)^{\gamma/(\gamma-1)}
 =\left(\frac{2}{\gamma+1}\right)^{\gamma/(\gamma-1)}.
\end{align*}
The pressure relation is isentropic. Its ratio is the
\emph{throat static pressure to inlet stagnation pressure};
it is not a universal downstream-to-chamber pressure test for
every valve or converging--diverging nozzle.

Steady continuity across the throat and the ideal-gas law give
\begin{align}
 \dot m_{\mathrm{out}}
 &=A_t\rho_tu_t
  =A_t\frac{p_t}{R_sT_t}\sqrt{\gamma R_sT_t}
  =A_t p_t\sqrt{\frac{\gamma}{R_sT_t}},\notag\\
 &=\frac{p_c A_t}{\sqrt{R_sT_c}}\sqrt{\gamma}
   \left(\frac{2}{\gamma+1}\right)^{
                  \frac{\gamma}{\gamma-1}-\frac12},\notag\\
 &=\frac{p_c A_t}{\sqrt{R_sT_c}}\sqrt{\gamma}
   \left(\frac{2}{\gamma+1}\right)^{
                  \frac{\gamma+1}{2(\gamma-1)}}.
   \label{eq:sb-throat-mass-derivation}
\end{align}
The two factors in the exponent come respectively from substituting
$p_t$ and $1/\sqrt{T_t}$. Define
\[
 \Gamma\equiv\sqrt{\gamma}
    \left(\frac{2}{\gamma+1}\right)^{\frac{\gamma+1}{2(\gamma-1)}},
 \qquad
 c^*\equiv\frac{\sqrt{R_sT_c}}{\Gamma}.
\]
Then, without suppressing the time-dependent throat area,
\begin{equation}\label{eq:sb-nozzle-flow}
 \dot m_{\mathrm{out}}(t)=\frac{p_c(t)A_t(t)}{c^*},
 \qquad R_sT_c=(\Gamma c^*)^2 .
\end{equation}
This reproduces Hill--Peterson Eq.~(3.14) \cite{hillpeterson}
and Sutton Eqs.~3--24 and 3--32 \cite{sutton}.
Under B2, this ideal $c^*$ is constant.

Here the chemical-composition assumption is frozen flow:
for a steady passage without species diffusion,
$Y_j(x)=Y_j(x_{\mathrm{in}})$ for each species mass fraction $Y_j$.
In a shifting-equilibrium model, composition instead adjusts to the
local thermodynamic state as the gas expands. Neither assumption
alone establishes $\mathfrak M_t=1$; sonic flow and chemistry must be specified
separately. A frozen mixture can still have temperature-dependent
heat capacities, so the constant-$\gamma$ formula above is a further
specialization. The algebraic definition of $c^*$ can also report a
measured or more detailed model's result; its numerical value must
come from that model, not automatically from this formula.

If the outlet is unchoked, the step $u_t=a_{\mathrm{s},t}$ is unavailable.
The mass flow then generally depends on downstream pressure as well
as upstream state and geometry. If the passage cannot adjust on the
imposed time scale, it needs an unsteady-flow model. In either case
retain $\dot m_{\mathrm{out}}$ in the chamber balance until the
appropriate outlet model has supplied it.

\paragraph{Pressure evolution before and after specifying the outlet.}
First use only the lumped density relation in B2 and the mass balance
in Section~\ref{sec:sb-mass-balance}:
\begin{align}
 \frac{dm_c}{dt}
 &=V_c\frac{d\rho_g}{dt}+\rho_g\frac{dV_c}{dt}
  =\dot m_g-\dot m_{\mathrm{out}},\notag\\
 \frac{V_c}{R_sT_c}\frac{dp_c}{dt}
 &=\dot m_g-\dot m_{\mathrm{out}}
      -\rho_g\frac{dV_c}{dt},\notag\\
 \frac{dp_c}{dt}
 &=\frac{R_sT_c}{V_c}
   \left(\dot m_g-\dot m_{\mathrm{out}}
                     -\rho_g\frac{dV_c}{dt}\right).
   \label{eq:sb-general-pressure-outlet}
\end{align}
For the present rigid grain, insert
$\dot m_g=\rho_pA_b r_b$ and $dV_c/dt=A_b r_b$:
\[
 \frac{dp_c}{dt}
 =\frac{R_sT_c}{V_c}
   \left[(\rho_p-\rho_g)A_b r_b-\dot m_{\mathrm{out}}\right].
\]
These equations require fixed $T_c,R_s$ as in B2, but no nozzle
assumption B3. If $N$ valves discharge directly from this chamber,
substitute $\dot m_{\mathrm{out}}=\sum_i\dot m_{\mathrm{out},i}$ here.
If they discharge from a receiving chamber, use that chamber's
inventory and thermodynamic state instead. Unknown outlet rates
still require additional relations or data.

Only for the choked nozzle just derived, substitute
Eq.~\eqref{eq:sb-nozzle-flow} and the burning law
Eq.~\eqref{eq:sb-burn-law} to obtain
\begin{align}
 \frac{dp_c}{dt}&=\frac{R_sT_c}{V_c(y)}
   \left[(\rho_p-\rho_g)A_b(y)\,a\,p_c^n-\frac{p_cA_t(t)}{c^*}\right],
   \label{eq:sb-pressure-ode}\\
 \frac{dy}{dt}&=a\,p_c^n .\label{eq:sb-web-ode}
\end{align}
Equation~\eqref{eq:sb-pressure-ode} is Hill--Peterson Eq.~(12.28), with
\eqref{eq:sb-nozzle-flow} used for the nozzle term. The $(\rho_p-\rho_g)$
factor comes from expanding the derivative in Sutton's Eq.~12--3:
part of the gas generated fills the volume uncovered by recession.
The variable $A_t(t)$ extends the fixed-area textbook specialization.
No $dA_t/dt$ term arises merely from substituting the instantaneous outlet
flow: $A_t(t)$ is in a flux, not in the stored inventory. If moving hardware
also changes the chamber volume, use the full $dV_c/dt$ including that
motion; then $V_c=V_c(y)$ and Eq.~\eqref{eq:sb-volume-rate} need revision.

The pair \eqref{eq:sb-pressure-ode}--\eqref{eq:sb-web-ode} is the complete
lumped model under B1--B5, with prescribed outlet area and fixed gas
temperature. It is not a complete ignition, energy, actuator or receiving-
chamber model. For a nozzle, thrust follows from
$F(t)=C_F(t)\,p_c(t)\,A_t(t)$, with $C_F$ evaluated for the current
geometry and pressure ratio; it need not remain constant as area changes.
A valve mass-flow law alone does not determine thrust. Burnout is
$y=w$; the tail-off after it requires $A_b(y)$ beyond $w$ (slivers) or a
blowdown model, and is outside this topic.

\subsection{Frozen-parameter equilibrium and a moving quasi-steady target}
\begin{proposition}[equilibrium pressure]\label{prop:sb-equilibrium}
Hold $y$ and $A_t>0$ fixed, assume B1--B5, and neglect $\rho_g$ against $\rho_p$. If
$n\ne1$, Eq.~\eqref{eq:sb-pressure-ode} has exactly one positive
equilibrium,
\begin{equation}\label{eq:sb-peq}
 p_{c,\mathrm{eq}}=\left(K a\rho_pc^*\right)^{\frac{1}{1-n}},\qquad K=\frac{A_b}{A_t}.
\end{equation}
If $n=1$, there is no isolated equilibrium: either every $p_c>0$ is one
(when $Ka\rho_pc^*=1$) or none is.
\end{proposition}
\begin{proof}
With $\rho_g$ neglected, $dp_c/dt=0$ is
$\rho_pA_bap_c^n=p_cA_t/c^*$. Divide by $A_t$ and by $p_c^n>0$:
\[
 Ka\rho_pc^*=p_c^{1-n}.
\]
If $n=1$ the right side is $1$ for every $p_c$, which gives the second
statement. If $n\ne1$, the map $p_c\mapsto p_c^{1-n}$ is a bijection of
$(0,\infty)$ onto itself, and raising both sides to the power $1/(1-n)$ gives
\eqref{eq:sb-peq}.
\end{proof}
Equation~\eqref{eq:sb-peq} is Sutton Eq.~12--6 solved for $p_1$, and
Hill--Peterson Eq.~(12.29) with their $(\rho_p-\rho_0)$ reduced to
$\rho_p$. Substituting $r_b=ap_c^n$ back into $Ka\rho_pc^*=p_c^{1-n}$ gives
$p_c=K\rho_p r_bc^*$, Sutton Eq.~12--4 in this frozen approximation.

\paragraph{Huzel and Huang's gas-generator pressure, Eq.~(4--47).}
Huzel and Huang, Section~4.6, printed p.~117 / downloaded PDF p.~134,
give the burning correlation and generation rate as Eqs.~(4--45) and
(4--46) \cite{huzelhuang}:
\begin{align}
 R_H&=k_1(T_b)\left(\frac{P_c}{P_{\mathrm{ref}}}\right)^n,
 &P_{\mathrm{ref}}&=1000\ \mathrm{psia},
 \label{eq:sb-huzel-burn}\\
 \dot w_g&=\rho_H A_bR_H.
 \label{eq:sb-huzel-generation}
\end{align}
The subscript $H$ distinguishes their recession speed $R_H$ from our
\emph{specific gas constant} $R_s$; it also distinguishes their propellant
weight density $\rho_H$ from our SI mass density $\rho_p$. Their printed
symbol is simply $R$ or $\rho$, without the editorial subscript.
The pressure ratio in \eqref{eq:sb-huzel-burn} is dimensionless:
$k_1$ is the measured recession speed at $P_{\mathrm{ref}}$ and the
specified grain temperature, not a universal propellant constant.
At fixed $T_b$ and in consistent SI units our Eq.~\eqref{eq:sb-burn-law}
has $a(T_b)=k_1(T_b)/P_{\mathrm{ref}}^n$.

The notation needed for the two cited equations is:
\begin{center}
\begin{tabular}{lll}
Huzel--Huang & Ours & Parallel notation\\ \hline
$R_H$ & $r_b$ & Sutton, Hill--Peterson: $r$; Humble: $r_b$\\
$P_c$ & $p_c$ & Sutton: $p_1$; Hill--Peterson: $p_0$\\
$A_0$ & $A_t$ & Hill--Peterson: $A^*$; Humble: $A_t$\\
$\rho_H$ & $\rho_p$ & Sutton: $\rho_b$; Humble: $\rho_p$\\
$\dot w_g$ & $\dot m_g$ & Hill--Peterson: $\dot m_g$\\
$k_1/P_{\mathrm{ref}}^n$ & $a(T_b)$ & Sutton, Hill--Peterson: $a$\\
\end{tabular}
\end{center}
The first and fourth rows identify the same physical quantities after
unit conversion. Huzel and Huang use $A_0$ for the \emph{outlet orifice},
not for the burning area. Their $c^*$ is our $c^*$. Their $\dot w_g$ is
a weight rate in lb/s, whereas our $\dot m_g$ is a mass rate in kg/s;
their Eq.~(1--32) uses $g$ to convert between the two conventions.
Williams' $m$ in Chapter~7 is a local burning \emph{flux}, not either
total rate; Turns does not assign a solid-grain counterpart here.

To recall Eq.~(1--32), begin with the sonic-throat calculation above,
Eq.~\eqref{eq:sb-throat-mass-derivation}. Under B3 it gives
\begin{align}
 \dot m_{\mathrm{out}}
 &=A_t\rho_tu_t
  =\frac{p_cA_t}{\sqrt{R_sT_c}}\,
    \underbrace{\sqrt{\gamma}
      \left(\frac{2}{\gamma+1}\right)^{
      \frac{\gamma+1}{2(\gamma-1)}}}_{\Gamma},
 \notag\\
 c^*&\equiv\frac{\sqrt{R_sT_c}}{\Gamma}
  =\frac{p_cA_t}{\dot m_{\mathrm{out}}}.
 \label{eq:sb-huzel-cstar-mass}
\end{align}
The first equality uses $u_t=\sqrt{\gamma R_sT_t}$, the ideal-gas law,
$T_t/T_c=2/(\gamma+1)$ and the isentropic throat pressure ratio, all
derived above. In Huzel and Huang's English weight-flow convention,
$\dot W_{tc}=g\dot m_{tc}$ when $\dot m_{tc}$ is measured in slugs/s.
Replacing the mass rate in \eqref{eq:sb-huzel-cstar-mass} by this weight
rate gives their Eq.~(1--32), printed p.~13 / downloaded PDF p.~32:
\begin{equation}\label{eq:sb-huzel-cstar-weight}
 c^*=\frac{(p_c)_{ns}A_tg}{\dot W_{tc}}
 \quad\Longleftrightarrow\quad
 \dot W_{tc}=\frac{(p_c)_{ns}A_tg}{c^*}.
\end{equation}
Here $(p_c)_{ns}$ is the stagnation pressure at the nozzle inlet;
identifying it with the gas-generator's $P_c$ neglects a pressure loss
between that chamber and its sonic orifice. We have used SI mass flow
in \eqref{eq:sb-nozzle-flow}, so its $g$ factor does not appear there.

Now specialize Eq.~\eqref{eq:sb-huzel-cstar-weight} to a choked orifice
of area $A_0$. For an equilibrium or a sufficiently slow quasi-steady
change, neglect chamber accumulation in
Eq.~\eqref{eq:sb-chamber-continuity}; only \emph{then} set generated
weight flow equal to outlet weight flow. Substituting
Eqs.~\eqref{eq:sb-huzel-burn} and \eqref{eq:sb-huzel-generation}
one at a time gives
\begin{align}
 \dot w_g&=\dot W_{\mathrm{out}},\notag\\
 \rho_HA_bR_H&=\frac{P_cA_0g}{c^*},\notag\\
 \rho_HA_bk_1\left(\frac{P_c}{P_{\mathrm{ref}}}\right)^n
   &=\frac{P_cA_0g}{c^*},\notag\\
 P_c^{\,1-n}
   &=\frac{\rho_Hk_1c^*}{gP_{\mathrm{ref}}^n}
     \frac{A_b}{A_0}.
 \label{eq:sb-huzel-pressure-power}
\end{align}
For $P_c>0$ and $n\ne1$, define the temperature-dependent constant
\begin{equation}\label{eq:sb-huzel-k2}
 k_2\equiv
 \left(\frac{\rho_Hk_1c^*}{gP_{\mathrm{ref}}^n}\right)^{1/(1-n)}.
\end{equation}
Raising \eqref{eq:sb-huzel-pressure-power} to $1/(1-n)$ yields precisely
their Eq.~(4--47):
\begin{equation}\label{eq:sb-huzel-pressure}
 \boxed{P_c=k_2\left(\frac{A_b}{A_0}\right)^{1/(1-n)}.}
\end{equation}
The units of $k_2$ are pressure, and $A_b/A_0$ is dimensionless.
In our notation the same result is Eq.~\eqref{eq:sb-peq}, with
$A_0\leftrightarrow A_t$ and $a=k_1/P_{\mathrm{ref}}^n$ after unit
conversion. For $n=1$, the preceding equation either has no positive
solution or permits every positive pressure at one special coefficient;
division by $1-n$ would be invalid.

Equation~\eqref{eq:sb-huzel-pressure} assumes a pressure-uniform chamber,
one choked outlet with fixed effective area, negligible accumulation,
and approximately fixed $k_1,n,\rho_H,c^*$ at the specified grain
temperature. Choking removes downstream pressure from the ideal outlet
flow law only while the critical pressure ratio is satisfied.
Constant $A_b$ alone does not ensure constant flow during ignition,
tail-off, or changes in orifice area, gas properties or grain
temperature. The transient pressure instead obeys the balance in
Eq.~\eqref{eq:sb-pressure-ode}; for a separate downstream plenum use
Eq.~\eqref{eq:sb-two-chamber-mass}.

When geometry or throat area varies, define the instantaneous target
\begin{equation}\label{eq:sb-moving-target}
 p_{c,\mathrm{eq}}(y,t)
 =\left[\frac{A_b(y)}{A_t(t)}\,a\rho_p c^*\right]^{1/(1-n)}.
\end{equation}
The transient pressure need not equal this target. Holding the other
parameters fixed, its logarithmic derivative along $y(t)$ is
\[
 \frac{d\ln p_{c,\mathrm{eq}}}{dt}
 =\frac{1}{1-n}\left(\frac{1}{A_b}\frac{dA_b}{dt}
                    -\frac{1}{A_t}\frac{dA_t}{dt}\right).
\]
This explicitly includes the changing outlet area. A valve on a separate
receiving chamber requires the coupled inventories in
Eq.~\eqref{eq:sb-two-chamber-mass}, not direct substitution of that valve's
area into a single motor-nozzle equilibrium formula.

\paragraph{The neglect of $\rho_g$.} With $p_c=7\ \mathrm{MPa}$,
$c^*=1550\ \mathrm{m/s}$ and $\gamma=1.2$, one has $\Gamma=0.6485$ and
$R_sT_c=(\Gamma c^*)^2=1.010\times10^6\ \mathrm{m^2/s^2}$, so
$\rho_g=p_c/(R_sT_c)=6.93\ \mathrm{kg/m^3}$. Against $\rho_p=1800\ \mathrm{kg/m^3}$
the ratio is $3.8\times10^{-3}$: the correction to $p_{c,\mathrm{eq}}$ is
a factor $(1-\rho_g/\rho_p)^{1/(1-n)}$, about $0.6\%$ at $n=0.35$. These inputs
are illustrative, not a particular propellant.

\paragraph{Sensitivity to the burning area.}
Take logarithms of \eqref{eq:sb-peq}:
$\ln p_{c,\mathrm{eq}}=(1-n)^{-1}(\ln K+\ln a+\ln\rho_p+\ln c^*)$. Holding
$a$, $\rho_p$, $c^*$ fixed,
\[
 \frac{\partial\ln p_{c,\mathrm{eq}}}{\partial\ln K}=\frac{1}{1-n}.
\]
With $A_t$ fixed, a grain whose burning area grows by 10\% raises the pressure by the factor
$1.1^{1/(1-n)}$: for $n=0.35$, $1.1^{1.538}=1.158$, a 15.8\% rise.
This amplification, not the burning law itself, is why grain geometry
(Sutton Section~12.3) is the dominant design variable of a solid motor.

\paragraph{Sensitivity to grain temperature.}
Hold $K$ fixed and differentiate the same logarithm with respect to $T_b$.
Only $a$ depends on $T_b$ (B4; Sutton notes $c^*$ changes only slightly,
p.~447, and we hold it fixed):
\[
 \pi_K\equiv\left(\frac{\partial\ln p_{c,\mathrm{eq}}}{\partial T_b}\right)_K
 =\frac{1}{1-n}\frac{d\ln a}{dT_b}=\frac{\sigma_p}{1-n},
 \qquad
 \sigma_p\equiv\left(\frac{\partial\ln r_b}{\partial T_b}\right)_{p_c}=\frac{d\ln a}{dT_b}.
\]
This is Sutton Eq.~12--14 and Hill--Peterson's $\Pi_p=\Pi_r/(1-n)$
(p.~602). With $\sigma_p=0.002\ \mathrm{K^{-1}}$ and $n=0.35$,
$\pi_K=0.0031\ \mathrm{K^{-1}}$: a motor fired at $-30^\circ$C instead of
$+20^\circ$C loses about $1-e^{-50\pi_K}\approx14\%$ of its chamber pressure.

\subsection{Frozen-parameter stability and the limits of quasi-steady tracking}
\begin{proposition}[pressure stability]\label{prop:sb-stability}
Under the hypotheses of Proposition~\ref{prop:sb-equilibrium}, with $y$
(hence $A_b$ and $V_c$) and $A_t$ frozen, write
$f(p_c)=\rho_pA_bap_c^n-p_cA_t/c^*$, so that
$dp_c/dt=(R_sT_c/V_c)f(p_c)$.
\begin{enumerate}
\item If $n<1$, every solution with $p_c(0)>0$ converges monotonically to
$p_{c,\mathrm{eq}}$, and the linearization there decays at the rate
$1/t_{\mathrm{rel}}$ with
\begin{equation}\label{eq:sb-tau}
 t_{\mathrm{rel}}=\frac{V_c c^*}{(1-n)R_sT_cA_t}=\frac{L^*}{(1-n)\,\Gamma^2c^*},
 \qquad L^*\equiv\frac{V_c}{A_t}.
\end{equation}
\item If $n>1$, $p_{c,\mathrm{eq}}$ is unstable: solutions starting below it
decay toward zero (extinction) and those above it grow without bound.
\end{enumerate}
\end{proposition}
\begin{proof}
Factor $p_c>0$ out of $f$ and use $\rho_pA_ba=A_t\,p_{c,\mathrm{eq}}^{1-n}/c^*$,
which is \eqref{eq:sb-peq} rearranged:
\[
 f(p_c)=p_c\left(\rho_pA_bap_c^{n-1}-\frac{A_t}{c^*}\right)
 =\frac{A_t\,p_c}{c^*}\left[\left(\frac{p_c}{p_{c,\mathrm{eq}}}\right)^{n-1}-1\right].
\]
If $n<1$, $(p_c/p_{c,\mathrm{eq}})^{n-1}>1$ for $p_c<p_{c,\mathrm{eq}}$ and $<1$
for $p_c>p_{c,\mathrm{eq}}$, so $f>0$ below and $f<0$ above. A scalar
autonomous equation cannot cross an equilibrium, so each solution moves
monotonically toward $p_{c,\mathrm{eq}}$ and, being bounded and monotone,
converges to an equilibrium, which can only be $p_{c,\mathrm{eq}}$. If $n>1$ the
signs reverse, which gives the second statement.

For the rate, differentiate $f$ at fixed $y$ and $A_t$, and evaluate at
$p_{c,\mathrm{eq}}$:
\[
 f'(p_{c,\mathrm{eq}})=n\rho_pA_bap_{c,\mathrm{eq}}^{n-1}-\frac{A_t}{c^*}
 =n\frac{A_t}{c^*}-\frac{A_t}{c^*}=-(1-n)\frac{A_t}{c^*},
\]
using $\rho_pA_bap_{c,\mathrm{eq}}^{n-1}=A_t/c^*$ again. With
$\delta=p_c-p_{c,\mathrm{eq}}$, the linearization is
$d\delta/dt=(R_sT_c/V_c)f'(p_{c,\mathrm{eq}})\,\delta=-\delta/t_{\mathrm{rel}}$ with the first
form of \eqref{eq:sb-tau}. The second form substitutes
$R_sT_c=(\Gamma c^*)^2$ from \eqref{eq:sb-nozzle-flow} and cancels one $c^*$.
\end{proof}
This is the argument Hill and Peterson draw as Fig.~12.18 (p.~601) and Sutton
states in one line on p.~446: $n<1$ for stability. The proposition adds the
global statement and the rate.

\paragraph{When a moving equilibrium can be followed.}
The proposition is about fixed parameters. For a slowly changing
$A_b(y(t))$ and $A_t(t)$, let
$\delta(t)=p_c(t)-p_{c,\mathrm{eq}}(y(t),t)$.
In the same small-$\rho_g/\rho_p$ model, local linearization gives
\[
 \frac{d\delta}{dt}
 \simeq-\frac{\delta}{t_{\mathrm{rel}}(t)}
       -\frac{dp_{c,\mathrm{eq}}}{dt},\qquad
 t_{\mathrm{rel}}(t)=\frac{V_c(y(t))c^*}{(1-n)R_sT_cA_t(t)},\quad n<1.
\]
The extra forcing term is the motion of the equilibrium itself.
Near the target, a useful local slow-change condition is
\begin{equation}\label{eq:sb-tracking-condition}
 t_{\mathrm{rel}}(t)\left|\frac{d\ln p_{c,\mathrm{eq}}}{dt}\right|\ll1.
\end{equation}
The coefficients must also vary slowly over the local response time.
This is a tracking estimate after transients have decayed, not a
global guarantee for arbitrary valve commands; directly check
Eq.~\eqref{eq:sb-storage-test} before discarding accumulation.

For illustration, $L^*=20\ \mathrm m$, $n=0.35$, $\gamma=1.2$ and
$c^*=1550\ \mathrm{m/s}$ give
$t_{\mathrm{rel}}=20/(0.65\times0.4206\times1550)=0.047\ \mathrm s$.
Changes much slower than this local response may admit a quasi-steady
approximation under the other assumptions. A long total burn does
not make a rapid valve movement quasi-steady. Since
$t_{\mathrm{rel}}\propto L^*=V_c/A_t$, a larger $V_c/A_t$ lengthens the response
time; a smaller ratio shortens it when the other parameters are fixed.
Ignition, tail-off, and changes on comparable or shorter time scales
require retaining the transient balance and checking the applicable
flow and thermal models. The same $L^*$ is the characteristic
chamber length used for liquid engines (Sutton Section~8.2).

\subsection{The first spatial correction: pressure drop along the port}
Relax B1 for a single cylindrical port of constant area $A$ and
circumference $c_p$ (Hill and Peterson write $c$), and keep the gas
isothermal at $T$. Here $p_c(x)$ and $\rho_g(x)$ are local static
gas pressure and density in the chamber port, rather than spatially
uniform chamber values. Let $p_h=p_c(0)$ be the head-end pressure
(Hill--Peterson's $p_1$); $\dot m=\dot m(x)$ is the local axial flow,
not the total surface-generation rate $\dot m_g$. Momentum for a slice of the port with no wall friction
is Hill--Peterson Eq.~(12.33), written in our notation as
$-A\,dp_c=d(\dot m u)$. Integrate from the head
end $x=0$, where $u=0$ and $p_c=p_h$, to $x$; with $A$ constant this gives
$p_h-p_c=\dot mu/A$ (Eq.~12.35). Substitute $u=\dot m/(\rho_g A)$ and
$\rho_g=p_c/(R_sT)$:
\[
 p_h-p_c=\frac{R_sT}{p_c}\left(\frac{\dot m}{A}\right)^2
 \quad\Longrightarrow\quad
 p_c^2-p_h p_c+R_sT\left(\frac{\dot m}{A}\right)^2=0 .
\]
With uniform burning, $\dot m=c_p r_b\rho_px$. The root that reduces to $p_h$
at $x=0$ is
\begin{equation}\label{eq:sb-port}
 \frac{p_c}{p_h}=\frac12\left[1+\sqrt{1-4R_sT\left(\frac{c_p r_b\rho_px}{p_hA}\right)^2}\,\right],
\end{equation}
Hill--Peterson Eq.~(12.37). The square root is real only while
$\dot m/A\le p_h/(2\sqrt{R_sT})$. At equality $p_c=p_h/2$ and
$u=\dot m/(\rho_g A)=\sqrt{R_sT}$, the isothermal sound speed: the port chokes
in this model at $\mathfrak M=1/\sqrt\gamma$. That limit, together with erosive
burning (Sutton Eq.~12--17; Hill--Peterson p.~603), is what forces a booster
model beyond zero dimensions: a long, slender grain early in the burn has a
high port Mach number, uneven recession along its length, and a head-end
pressure above the aft-end pressure.

\subsection{How much chemistry does a solid-motor simulation carry?}
In the model above the chemistry appears in exactly two places: $c^*$ (and
$\gamma$, $C_F$), which is a thermochemical equilibrium result, and the pair
$(a,n)$, which is measured. The honest answer to ``how high-fidelity is
solid-propellant chemistry'' is set by the second.

\paragraph{What the books say.} Sutton, p.~445--446: the burning-rate
relations are ``empirical generalizations'', and ``strict analytical modeling
and supportive research have yet to adequately predict the burning rate of a
new propellant in a new rocket motor.'' Rates are measured in strand
burners, then in small ballistic evaluation motors, and finally in full-scale
firings; strand rates run 4 to 12\% below full-scale, and the relation
between the three ``must be determined empirically for each propellant
category and grain configuration'' (p.~444). Hill and Peterson likewise call
$r_b$ ``an empirically determined function'' (p.~598).

\paragraph{Why the chemistry is hard to carry.} Williams Chapter~7 is the
theory. A homogeneous propellant deflagrates through a condensed-phase
reaction zone, a surface where gasification occurs, and a gas-phase flame
(Sections~7.1--7.5); Williams treats the regimes in which the condensed
phase or the gas phase controls the rate. A \emph{composite} propellant
(Section~7.7) is not homogeneous at all: oxidizer crystals (ammonium
perchlorate, tens to hundreds of micrometres) sit in a polymer binder, so the
burning surface carries a premixed oxidizer flame and diffusion flames
between oxidizer and binder vapours, with a length scale set by particle
size. Three facts follow.
\begin{enumerate}
\item The \emph{gas-phase} flames can be given detailed elementary
kinetics, as in a liquid-engine flame, and research models do this for the
oxidizer's own flame. The \emph{condensed-phase} decomposition and the
surface are represented by a few global steps, because their elementary
chemistry is not known to comparable precision.
\item The heterogeneous geometry must be modelled statistically (particle
size distribution, packing) before any kinetics apply. The classic
engineering form is the multiple-flame model of Beckstead, Derr and Price
\cite{BDP1970}; modern research resolves packed oxidizer particles in
three dimensions with reduced kinetics.
\item The flame zone is micrometres to a fraction of a millimetre thick,
against a motor metres long. No full-motor simulation resolves it.
Full-motor CFD imposes the burning rate as a surface boundary condition,
$r_b=r_b(p_c,\text{local flow},T_b)$, calibrated as above.
\end{enumerate}

Aluminized propellants add a second, largely empirical layer: aluminium
particles agglomerate on the surface and burn as droplets in the port, and the
alumina smoke they leave creates the two-phase nozzle loss (Hill--Peterson
Section~12.6, p.~598; Sutton Section~3.5) and damps acoustic modes (Williams
Section~9.1.4.7). Turns Chapter~14 (burning of carbon, one-film and two-film
models) \cite{turns} is the closest textbook analogue for a burning metal particle.

\paragraph{The fidelity ladder, solid versus liquid.}
Ordered by what is resolved rather than assumed:
\begin{enumerate}
\item \emph{Equilibrium thermochemistry} for $T_c$, $\gamma$, $c^*$,
$C_F$, with frozen or shifting expansion (Sutton Chapter~5; the NASA CEA
program \cite{CEA1994} is the standard reference implementation). Same for
both.
\item \emph{Lumped internal ballistics}, Eqs.~\eqref{eq:sb-pressure-ode}--%
\eqref{eq:sb-web-ode}, with a grain-geometry model for $A_b(y)$, $V_c(y)$.
The liquid analogue is the lumped engine network (tanks, lines, pumps,
chamber) of Huzel and Huang Section~10.2 \cite{huzelhuang}.
\item \emph{Quasi-one-dimensional port flow} with erosive burning, axial
pressure drop and local recession, Eq.~\eqref{eq:sb-port} and its
generalizations. This is the working fidelity of solid-motor design codes.
\item \emph{Finite-rate nozzle chemistry and two-phase flow}: gas-particle
drag and heat-transfer lag in the nozzle (solid), recombination kinetics
in the expansion (both; Williams Section~4.2, Hill--Peterson
Section~12.4).
\item \emph{Multidimensional chamber flow and acoustics}: port vortex
shedding, combustion response functions, instability (Williams
Section~9.1; Sutton Section~14.4). For liquids, injector-resolved CFD with
tabulated or reduced chemistry.
\item \emph{Resolved flames}: the level at which each elementary reaction
rate is computed. For liquids this is reachable in canonical problems
(well-stirred and plug-flow reactors, counterflow flamelets; Turns
Chapter~6) and, at high cost, in combustor LES. For solids it exists only in
research sandwich and packed-particle simulations, and even there the
condensed phase remains global.
\end{enumerate}
So the asymmetry the question points at is real, with one correction: in
both cases a full-engine simulation does not compute every elementary rate
everywhere. Liquid-engine chemistry is detailed \emph{where it is used}
(flamelet tables, reduced mechanisms, finite-rate nozzles), because the
propellants are gases or sprays of known molecules. Solid-motor chemistry
is empirical \emph{at the burning surface}, because that surface is a
heterogeneous, partly condensed-phase reaction zone whose rate is still
measured rather than predicted.

\paragraph{Scope of this topic.} Propositions~\ref{prop:sb-equilibrium} and
\ref{prop:sb-stability} are proofs under B1--B5 at frozen grain geometry and throat area. The
numerical values are illustrative inputs chosen here, not data for a named
propellant; no executable check has been run for this topic. The
descriptions of research practice in the last subsection are a summary
pointing to the literature named, not results derived in these notes.
